\section{Strip Partitioning and Mapping}

% JW (acted, archived 2026-07-22 -- retitled, and the architecture
% figure sits with this section):
% \JW{title should be Tessera Architecture, move figure 2 here to this section}
\label{sec:design}



\figreal{3tesseraarch}{arch}{1.0}{%
\textsc{Tessera} architecture and weight mappings at $D{=}4$, $d{=}2$.}{%
(a)~The placed array: strips of stacked $d{\times}d$ blocks,
flanked by strip memories, with twist and strip boundary PEs.
(b)~A square $4{\times}4$ weight maps onto all the blocks, the
strips merging as one array. (c)~A tall $4{\times}2$ and (d)~a wide
$2{\times}4$ weight matrix each map onto a two-block band.}



The fabric streams an orthogonal wavefront: weights and activations
enter at one edge, and results drain at the opposite edge. Each strip
therefore reads its activations and weights from its left edge and
writes its outputs at its right edge, so the wavefront control stays
simple and adjacent strips can share a strip memory. Because
everything moves in that one horizontal direction, the array is
partitioned along columns into vertical strips, with a strip memory
interleaved between adjacent strips to serve the activation, weight,
and output traffic of the strips it flanks. The chip lays the strips
out horizontally and backs them with a main SRAM buffer.

Each strip is a $D{\times}d$ block of PEs running the twist-torus
dataflow of \S\ref{sec:ttsa:torus}, built from $d{\times}d$ PE
blocks stacked along the column, and the buffers on its two sides
serve operand reads and partial-sum writes. Composing strips
exposes partition points in both directions, between adjacent strips and
between the stacked blocks. Each boundary carries one binary
choice, partition or merge, and the controller sets the choice per
boundary. The fold timing of \S\ref{sec:ttsa:anchors} applies per
square, with $D$ replaced by the square side.

The unit this section builds toward is the flex compute region
(FCR), a rectangle bounded by those cuts at $d$ granularity in
both directions, rows $n_1 d$ to $n_2 d$ of a contiguous strip
group. The twist boundary PE (TBPE) cuts the FCR's row edges. The
strip boundary PE (SBPE) decides whether data crosses to the next
strip for reuse or drains into the strip memory. Each region is
reserved for one workload between two reform points, the moments
the partitioning changes, and multiple regions co-exist and are
reformed between GEMMs.

\subsection{Strip memory}
\label{sec:design:mem}
To the array, each flanking strip memory is one logical bank,
shown in Fig.~\ref{fig:arch}a, and every strip memory connects
downward to the main buffer over the system interconnect.
Internally it holds two single-port R/W banks in double-buffered
ping-pong. At any time one bank serves the activation and weight
reads and the output writes, while the other fetches new data from the
main buffer, and the roles swap between tiled GEMMs. A merging boundary reuses activations and
partial sums in-array through the
boundary-enable muxes, so only outputs
need to reach the SRAM, and fresh activations feed the next strip
directly. Each bank is divided into $D/d$ lanes at the
partition granularity, one lane per $d$-row block, and lane $y$
wires directly to array rows $[yd, yd{+}d)$.
A lane holds $d$ operand words of $16$ bits or $d$ partial sums of
$32$ bits and serves them in one SRAM access. The array therefore
reads its $D$ operands and drains its $D$
partial sums each cycle from $D/d$ disjoint macros in parallel,
without a bank conflict. Aggregate bandwidth scales with the strip
count because it is distributed across these private lanes rather
than gated by a shared port.

The same ping-pong hides the reloads. Whichever bank the array is
not reading is filled or drained through an independent back port
to the main buffer. Refill therefore overlaps compute whenever a
tiled GEMM streams at least
as long as its reload, the same $M{\ge}d$ condition that hides a
fold's preload (\S\ref{sec:ttsa:anchors}), and a boundary-partitioned strip
drains its partial sums into the neighboring output lane rather than
off-chip. The banks are compiled SRAM
macros, and their per-lane glue is the memory-logic area and power
counted in \S\ref{sec:eval:tradeoff}.

% Placement: declared here so the [!t] float lands on the \S IV-B
% page (user 2026-07-29: Fig 5 on p5).
\figreal{4PEs}{pes}{1.0}{%
The PE micro-architecture and datapaths.}{%
(a)~TTSA PE: a fused multiply--add core plus a weight-bypass
register and a lock register (L). (b)~Strip boundary PE (SBPE):
enable (E) forwards the partial sum to the next strip or drains it
to strip memory.
(c),~(d)~Twist boundary PE (TBPE), odd/even-column forms.}

\subsection{Partitioning or merging: boundary PEs}
\label{sec:design:boundary}
\label{sec:design:cut}

% Placement: moved to the \S IV-B head (2026-07-31) so Fig 6/7 land one
% page earlier.  5Tallandwide = the weight loads, 6areuseandpreuse =
% the two reuse paths, both from fig/new/.
\figrealstar{5Tallandwide}{tallwide}{1.0}{%
Weight load for tall and wide weight matrix.}{%
(a)~Tall weight matrix ($K{=}8$, $N{=}4$): each strip memory holds one
depth-$d$ chunk of $B$, locked into its strip in the twist
permutation order. (b)~Wide weight matrix ($K{=}4$, $N{=}8$): the two
width-$d$ chunks load the same way, one per strip.}

\figreal{6areuseandpreuse}{reuse}{1.0}{%
In-array reuse across a merging strip boundary.}{%
(a)~Partial-sum reuse for a tall weight matrix: partial sums cross the
boundary and accumulate to a single $C$ exit. (b)~Activation reuse for a wide weight matrix:
the activation forwards across the boundary, and each strip drains
its own $C$ columns.}

Fig.~\ref{fig:pes}b--d shows the two types of
boundary PEs that implement the partition-or-merge choice, one per
direction of an FCR, and the muxes and demuxes each
carries on the paths that cross the partition line. In each, a
boundary-enable register steers those paths, while the FMA, weight, and lock registers are those of an interior PE. The
controller sets both enables per region.

\paragraph{Twist boundary PE.}
The TBPE sits on the strip-internal activation chain at
every $d$-row block boundary along $D$, in an odd- and an even-column
form, drawn in Fig.~\ref{fig:pes}c,d. The twist's crossing links
alternate sides with column parity. The two forms therefore carry
their muxes on complementary path sets. The TBPE's choice sets a region's top
and bottom row at $d$ granularity. To merge, the TBPE's mux selects the
across-boundary activation, so that the blocks above and below the line
run as one taller region. To partition, the TBPE selects the in-strip predecessor
and restarts the chain, which opens a new region edge. Either choice
preserves the torus-link timing of \S\ref{sec:ttsa:torus}.

\paragraph{Strip boundary PE.}
The SBPE sits at the end of a PE row on the strip's
memory-facing edge. The SBPE chooses whether the dataflow crosses the
boundary for reuse or exchanges its data with the intervening
strip memory.
The SBPE's output fans out to a pipeline flop forwarding to the next strip when
merging, and to a direct write into the adjacent strip memory when
partitioned, as Fig.~\ref{fig:pes}b shows. When merging, the two strips run
as one wider array. When partitioned, the crossing links go inactive and each
strip runs independently, with the left strip draining into the strip
memory that serves the right strip's operand reads.

Either boundary therefore carries one binary decision, and the
same one-bit enable drives every mux along it. A partition at a strip
boundary is then a row of multiplexers,
each with an optional pipeline flop. The per-PE datapath, per-strip
accumulator,
and in-strip torus links are untouched. \textsc{Tessera}'s boundary cost
grows only with the number of boundaries, one mux row per strip
boundary and per block edge, and is independent of the array side. On
the classical mesh, the skew registers and per-PE multiplexers
replicated into every sub-array multiply as the granularity
becomes finer.

Fine granularity is therefore cheap in silicon, but each
boundary still carries a data-reuse decision. Partitioned regions no longer
share operands or chain partial sums in-array. A shared
activation is read once per region instead of once per array, and
partial sums that once accumulated across the boundary detour through
the strip memory, so energy per operation rises at finer
granularity (\S\ref{sec:eval:tradeoff} measures this balance).
Partitioning the array is therefore only part of the problem. A
mapping and schedule must still feed the regions and accumulate
each GEMM across them, and the rest of this section builds both.

\subsection{Mirror-chain: mapping non-square weight matrices}
\label{sec:design:mapping}

A weight matrix is the $K{\times}N$ weight ($B$) block of a tiled GEMM,
$K$ its reduction depth and $N$ its output width. The weight matrix is
\emph{wide} when $N{>}K$ and \emph{tall} when $K{>}N$, the square
$K{=}N$ being the special case. Non-square weight matrices are the common case
in serving. The projections and attention of an LLM are naturally
non-square, splitting a large GEMM in $D{\times}D$ units leaves
non-square edge remainders, and the
partition minimizing the energy--delay product (EDP) often lands
at a non-square granularity.

On a partitioned array, each weight matrix takes its own region of the
array. Prior work on this
dataflow~\cite{Abdelmaksoud_2026} pads every non-square weight to the
full $D{\times}D$ square, which lowers utilization. Cutting a non-square weight matrix into regions naively
breaks the reuse a classical systolic array would give it. A partition along $K$
stops the partial sums from flowing region to region, so every
region writes its partial sums to SRAM and the results are
re-accumulated. A partition along $N$ stops the activation from flowing,
so every region reads the same $A$ from SRAM again. \textsc{Tessera} is the
first to map non-square weight matrices onto this dataflow with the in-array
reuse intact.

Prior partitioned designs add hardware to adapt non-square shapes to
the diagonal wavefront, such as the omni-directional multiplexers of
Planaria~\cite{ghodrati2020planaria}. Our orthogonal wavefront instead
handles them in the mapping alone, with no added hardware.
The \emph{mirror-chain} mapping does this in two steps. A $K{\times}N$ weight matrix first pads
to the nearest $s{\times}l$ rectangle, where the short side $s$ is
a multiple of $d$ and the long side $l$ is a multiple of $s$, and
the padded weight matrix then further splits into $s{\times}s$ square blocks, laid horizontally
across the strips on a band of rows. The weight load runs in parallel
for each $s{\times}s$ block as the tiled weight matrix is split. As
Fig.~\ref{fig:tallwide} shows, the leading block loads in the twist
permutation order and the trailing block in the reversed order, so
their weight layouts are mirror-symmetric. This mirrored layout is what
lets neighboring $s{\times}s$ blocks reuse partial sums and activations
across the boundary. A split along $K$ gives a tall weight matrix,
whose blocks chain their partial sums. A split along $N$ gives a
wide weight matrix, whose blocks share the activation
(Fig.~\ref{fig:reuse}). The two cases give
the two mappings below, partial-sum reuse for tall weight matrices and
activation reuse for wide ones.

% chenyi9: addition start -- general mirror-chain definition; preserve submitted text.
\begingroup
\textcolor{red}{For an arbitrary chain, number the square weight blocks $B_b$ from
$b=0$. Let $J$ be the $s\times s$ reversal matrix, with $J^2=I$.
The weight matrix passed to the existing twist-order loader is}

\begin{equation}
\color{red}
 \widehat B_b=J^b B_b J^b.
 \label{eq:full-mirror}
\end{equation}

\textcolor{red}{Even blocks use regular order. Odd blocks reverse both local $K$ and
$N$ indices: $(JB_bJ)_{k,n}=(B_b)_{s-1-k,s-1-n}$.
Thus ``reverse order'' specifies the weight indices, rather than
only the direction in which a buffer is drawn.}

\textcolor{red}{Across an $s$-column square, including its outgoing edge, both routes
reverse physical row order. The alternating adjacent swaps move a
token toward an endpoint, pause it for one cycle, and reverse its
direction. This reflected walk has period $2s$; after $s$ edges,
row $r$ becomes $s-1-r$, for either starting phase and either parity
of $s$. Each connection therefore applies $J$, changing the order
from $J^b$ to $J^{b+1}$. The same alternation extends to any number
of connected blocks.}
\par
\endgroup
% chenyi9: addition end

\subsubsection{Tall weight matrices, partial-sum reuse}
\label{sec:design:mapping:klarge}
Figure~\ref{fig:tallwide}a shows a tall weight matrix, $K{=}8$, $N{=}4$, on a
$4$-row band across two strips of width $4$. The $K$ side splits
into two $4{\times}4$ blocks laid along the strips. The left strip holds
$B_{0..3}$ stationary and the right strip $B_{4..7}$, each preloaded
from its adjacent strip memory, the left strip in the regular twist
permutation order and the right strip in the reversed order. The activation also splits
along $K$ into an upper and a lower half, one per strip. Each
strip receives its own $K$-slice of the activation at its left
boundary, the second half $A_{04}$ to $A_{07}$ entering from the strip
memory at the strip boundary. Partial sums chain rightward across the
merging boundary, accumulating through the second strip to generate
the output matrix $C$ at the rightmost strip memory (Fig.~\ref{fig:reuse}a).

% chenyi9: addition start -- tall alignment identity; preserve submitted text.
\begingroup
\textcolor{red}{For the general tall case, let $A_b$ contain the activation columns
matching $B_b$, and let $S_b$ be the partial sum before block $b$,
with $S_0=0$. Incoming partial sums and fresh activations have orders
$S_bJ^b$ and $A_bJ^b$. Computation followed by the outgoing edge gives}

\begin{equation}
\color{red}
 \bigl[S_bJ^b+(A_bJ^b)\widehat B_b\bigr]J
   =(S_b+A_bB_b)J^{b+1}.
 \label{eq:tall-invariant}
\end{equation}

\textcolor{red}{The updated sum has exactly the next block's required order.
Each partial sum therefore continues to accumulate weights with the
same output-column index $n$, proving the chain invariant by induction.}
\par
\endgroup
% chenyi9: addition end

\subsubsection{Wide weight matrices, activation reuse}
\label{sec:design:mapping:ksmall}
Figure~\ref{fig:tallwide}b shows the transposed case, a wide weight matrix of
$K{=}4$, $N{=}8$ on the same band. The $N$ side splits into two
$4{\times}4$ blocks laid along the strips. The left strip holds
columns $0$--$3$ stationary and the right strip columns $4$--$7$, each
preloaded from its adjacent strip
memory, the left strip in the regular twist permutation order and the
right strip in the reversed order. The activation enters once at the left edge, crosses strip 0,
and is forwarded over the merging boundary. Every block
therefore reuses the same $A$ (Fig.~\ref{fig:reuse}b). The left strip drains output columns $C_0$ to $C_3$ and the right
strip columns $C_7$ down to $C_4$, strip 0 into the
strip memory at the strip boundary and strip 1 at its right edge.

% chenyi9: addition start -- wide alignment identity and comparison; preserve submitted text.
\begingroup
\textcolor{red}{In the wide case, every block uses the same activation matrix $A$.
Block $b$ receives $AJ^b$ and computes}

\begin{equation}
\color{red}
 (AJ^b)\widehat B_b=(AB_b)J^b.
 \label{eq:wide-invariant}
\end{equation}

\textcolor{red}{The outgoing activation becomes $AJ^{b+1}$, which is the next block's
required input. Here the $K$ mirror preserves the reduction-index
alignment, whereas in the tall case the $N$ mirror preserves the
partial-sum alignment. The wide case's $N$ mirror accounts for the
reversed output-column order. Regular weights in every block would
break these respective alignments.}
\par
\endgroup
% chenyi9: addition end

Either orientation retires useful MACs at the same operand bandwidth
as a $D{\times}D$ square reading the same shape padded to $D$, while the
fixed square fills the remaining cycles with padded zeros. Both
mappings advance at one MAC per PE per cycle, so a GEMM's run time
is set by the volume its mapping executes, the padded volume on the
fixed square against only the useful volume on the band. The
speed-up on an underutilized shape is therefore the ratio of padded
to useful MACs, the reciprocal of the $\min(M,N,K)/D$ utilization of
\S\ref{sec:bg:llm}, and $D$-aligned shapes perform equally.

\subsection{Mapping and scheduling}
\label{sec:design:batch}
\label{sec:design:modes}
\label{sec:design:algo}

A GEMM larger than the array is split the same way: with a
$K{\times}N$ weight matrix, $K$ splits
into $h$ blocks of size $k$ and $N$ into $w$ strips of size $n$,
with $k$ and $n$ dividing $D$, and each $k{\times}n$ weight matrix maps
onto an FCR through the two patterns of
\S\ref{sec:design:mapping}.

The choice of $(k,n)$ is made per GEMM by EDP, and a finer
granularity admits more candidates. Modern systolic accelerators
back the array with a large on-chip SRAM~\cite{jouppi2023tpu},
so the energy that moves with the choice is the SRAM traffic. A
smaller $n$ splits $N$ into more strips and rereads the activation
once per strip, so read energy grows with the strip count
$w{=}\lceil N/n\rceil$. A smaller $k$ splits $K$ into more blocks
and writes one partial sum per block, so write-and-reduction
energy grows with the block count $h{=}\lceil K/k\rceil$. For
example, with a large $M$ the choice leans to a larger $n$, saving
rereads of the $M$-row activation stream; with a small $M$ it
leans to a larger $k$, saving partial-sum writes and
re-accumulation.

The delay side, for a chosen array size, comes from the per-weight-matrix
fold time of a weight-stationary array~\cite{samajdar2020systematic}: under
the mirror-chain mapping a padded
$k{\times}n$ weight matrix of $H$ rows $\times$ $W$ columns, $H$ the
shorter edge, is a chain of $W/H$ square blocks and folds skew-free
in $H{+}M{+}W{-}1$ cycles. Across a GEMM's $R$ rounds only the
last exposes its traverse tail. Every round uses the same $(k,n)$, so
$H$ and $W$ stay constant, and
\[
  T_{\mathrm{GEMM}}
    \;=\; H \;+\; R\,\max(M, H) \;+\; W \;-\; 1 ,
\]
where each round costs $H$ when the fold tail cannot be hidden and $M$
when double buffering hides it ($M{\ge}H$).

% chenyi9: addition start -- timing interpretation; preserve submitted text.
\begingroup
\textcolor{red}{This expression reserves a full $\max(M,H)$ interval for each round.
When $M<H$, the last valid output precedes the end of the final
reserved interval by $H-M$ cycles. Extra boundary pipeline stages,
if present, add their latency to the traverse tail.}
\par
\endgroup
% chenyi9: addition end

% Placement: [!ht] so the float lands beside its prose (user placement
% rule 2026-07-22: keep it next to this subsection, not forced to a page
% top/bottom).
% Codex: float color guard start -- preserve the submitted algorithm color.
\begingroup
\makeatletter
\let\TesseraSavedFloatMakebox\float@makebox
\def\float@makebox#1{\color@vbox\normalcolor\TesseraSavedFloatMakebox{#1}\color@endbox}
\makeatother
\begin{algorithm}[t]
\footnotesize
\caption{Per-GEMM EDP-optimal mapping.}
\label{alg:tessera-sched}
\begin{algorithmic}[1]
\REQUIRE GEMMs $(M_\ell,K_\ell,N_\ell)_{\ell=1}^{L}$; array side $D$;
granularity $d$; sizes $\mathcal{S}{=}\{s : s{=}d\cdot 2^n,\ n{=}0,1,\ldots,\log_2(D/d)\}$
\ENSURE rounds $\mathcal{R}$ with per-weight-matrix placements
\FORALL{GEMMs $\ell$}
  \STATE $(k,n) \gets \arg\min_{(k,n)\in\mathcal{S}^2}\mathrm{EDP}_\ell(k,n)$
  \STATE split into $k{\times}n$ weight matrices:
         $h{=}\lceil K_\ell/k\rceil$ blocks $\times\ w{=}\lceil N_\ell/n\rceil$ strips
\ENDFOR
\STATE $r \gets$ empty $D{\times}D$ round
\FORALL{weight matrices $t$ in arrival order}
  \IF{$t$ fits a free $k{\times}n$ slot in $r$}
    \STATE place $t$ in $r$
  \ELSE
    \STATE close $r$;\ open a new $D{\times}D$ round;\ place $t$
  \ENDIF
\ENDFOR
\RETURN all rounds $\mathcal{R}$
\end{algorithmic}
\end{algorithm}
\endgroup
% Codex: float color guard end


Algorithm~\ref{alg:tessera-sched} picks per GEMM, over $\mathcal{S}^2$,
the $(k,n)$ minimizing the energy--delay product $\mathrm{EDP}$. Weight matrices are then first-fit into $D{\times}D$ rounds and packed
across GEMMs, and a partly-filled round stays open for the next
independent GEMM.

Finer granularity runs faster, as a smaller weight matrix fills more of the array
and $T$ falls, but it loses reuse and thus energy. \textsc{Tessera} can still go
finer than prior designs because on serving's small, non-square shapes the
EDP keeps improving at finer granularity, so the per-GEMM $\arg\min$
lands well below the full array. EDP stops shrinking the weight matrix once
the utilization recovered no longer outweighs the reuse lost.

\subsection{Scale-up}
\label{sec:scale}
\label{sec:scale:strip}
% fig 6 (4scaleupview, fig:scaleup, chip-scale view) deleted
% 2026-07-29 per author note; 4scaleupview.pdf stays in fig/.
Growing $D$ stacks more
$d{\times}d$ PE blocks inside each strip, growing the number of strips appends another
strip and strip-memory pair, and a strip's merging neighbor is always
adjacent. The critical path therefore stays at one pipeline stage and
the minimum partition size at $d$ for any $D$ and $n$. The two directions
scale independently. Because the EDP optimum needs only a short
strip, a chip can grow by replicating strips
horizontally rather than raising $D$, keeping the array short and
its physical implementation simple.